\section{Caída de costos, la paradoja de Jevons y el contrabando industrial}

\subsection{La caída de costos de 1200 veces y la paradoja de Jevons}

El costo de inferencia en la era de GPT-3 era de \$60-70 por millón de tokens. Tres años después (mediante modelos pequeños como Llama 3B), la misma capacidad bajó a unos cuantos centavos: una reducción de \textbf{1200 veces}. DeepSeek está en la línea de tendencia de costos de un modelo del nivel de GPT-4, pero fue el primero en llegar a ese punto de precio.

Las acciones de NVIDIA se desplomaron tras el lanzamiento de DeepSeek, pero eso obedeció a un contagio social basado en una narrativa equivocada. Ningún modelo público ha costado más de \$1000 millones en entrenamiento. La afirmación de que «DeepSeek solo costó 6 millones de dólares» es una mala lectura del artículo: eso es únicamente el costo de horas de GPU de una sola corrida de preentrenamiento, sin incluir investigación y desarrollo, experimentación, entrenamiento posterior ni despliegue de inferencia.

\begin{importantbox}{La paradoja de Jevons confirmada a la perfección en la IA}
La paradoja de Jevons en economía (planteada en 1865): cuando mejora la eficiencia con la que se usa un recurso, el consumo total aumenta en lugar de disminuir, porque el menor costo estimula más demanda.

La eficiencia de la máquina de vapor mejora $\to$ el consumo de carbón no baja, sube. El costo de inferencia de IA baja $\to$ la demanda de GPU no baja, sube.

Evidencia directa tras el lanzamiento de DeepSeek:
\begin{itemize}
    \item El precio de renta de H100 en AWS no bajó, \textbf{subió}
    \item El H200 casi se agotó
    \item Todos los proveedores de nube reportan una demanda disparada de GPU
    \item Nuevos casos de uso (agentes intensivos en inferencia, RAG a gran escala) se volvieron viables económicamente
\end{itemize}

\textit{``'Las Scaling Laws murieron' duró un mes... y luego salieron O1, O3, R1, y ahora se volvió 'los modelos avanzan demasiado rápido'.''} (Dylan)

Inferencia más barata $\to$ más casos de uso $\to$ más demanda de inferencia $\to$ se necesitan más GPU. No es un juego de suma cero: el mercado total se está expandiendo. Por eso las perspectivas de largo plazo de NVIDIA no se oscurecen por DeepSeek.
\end{importantbox}

\subsection{La escala y los métodos del contrabando de GPU}

El contrabando de GPU pasó de ser una zona gris a una operación de nivel industrial, con múltiples canales:

\textbf{Canal de renta en la nube}: ByteDance es el «contrabandista» más grande: renta más de 500,000 GPU a nivel mundial a través de Oracle, Google y diversos proveedores de nube pequeños. Reparte la operación entre entidades en varios países para evitar el escrutinio de una sola jurisdicción. Este «contrabando distribuido» es extremadamente difícil de rastrear porque cada contrato de renta individual queda por debajo del umbral regulatorio.

\textbf{Contrabando físico}: alguien viajó en primera clase desde el aeropuerto de San Francisco a Shanghái con cajas de servidores GPU de SuperMicro; esto no es ficción, ocurrió de verdad. Se estima que en 2024 entre 200,000 y 300,000 GPU llegaron a China a través de pequeñas empresas en Singapur y Malasia.

Los límites de las \textbf{nuevas reglas de difusión}: restringen a las entidades vinculadas con China a rentar menos de 2000 GPU o comprar menos de 1500 GPU. Pero el reto de aplicación es enorme: ¿cómo rastrear al usuario final de recursos de cómputo en la nube a escala global? ¿Cómo se define una «entidad vinculada con China»? ¿Una empresa registrada en Singapur pero operada por ciudadanos chinos cuenta o no?

\textit{``Las GPU van a superar a las drogas y a las armas como el bien de contrabando de mayor valor por kilogramo.''} (Dylan) Una H100 cuesta cerca de \$30,000 y pesa menos de 5 kg, más de \$6,000 por kilogramo, muy por encima de la mayoría de los bienes prohibidos. Y a diferencia de las drogas, las GPU son mercancía legal en algunos mercados, solo restringida en ciertas direcciones de exportación. Eso crea un espacio de arbitraje singular.

\subsection{Destilación, espionaje y flujo de conocimiento}

\subsubsection{La naturaleza de la controversia sobre la destilación}

La destilación (entrenar un modelo más débil con las salidas de uno más fuerte) es una práctica estándar en la industria. OpenAI afirmó que DeepSeek usó las salidas de su API, y el Financial Times lo publicó en un titular. Pero Nathan señala que esto se parece más a un «control de la narrativa» que a una acusación técnica real:

\textit{``¿Por qué es poco ético que yo entrene con las salidas de tu modelo, mientras que tú puedes entrenar con el texto de internet?''} (Nathan) El contenido generado por usuarios a través de ChatGPT se copia y pega en internet, y así se vuelve un dato público. Muchas empresas emergentes estadounidenses entrenan abiertamente sobre salidas de OpenAI; ¿por qué ellas sí pueden y DeepSeek no?

Una pregunta más profunda: \textbf{¿a quién pertenece la propiedad intelectual de las salidas de un modelo?} Si yo genero un fragmento de código con GPT-4, ¿el derecho de autor de ese código es de OpenAI o mío? Esta pregunta todavía no tiene una respuesta definitiva en el terreno legal.

\subsubsection{Movilidad de talento y espionaje industrial}

La cultura de cambio de empleo de Silicon Valley (las cláusulas de no competencia son ilegales en California) hace que las ideas —no el código— circulen libremente. Un caso concreto: un ingeniero que ayudó a Google a construir el contexto de un millón de tokens de Gemini se cambió a Meta, y se espera que la siguiente generación de Llama tenga un contexto de un millón de tokens. ¿Esto es transferencia legítima de conocimiento o robo de tecnología de facto? El límite es sumamente difuso.

Robar código o datos es un delito y es relativamente fácil de rastrear. Pero \textbf{la transferencia de ideas} es inevitable: no se puede «olvidar» la intuición arquitectónica que se aprendió en la empresa anterior.

\begin{knowledgebox}{Japón: un paraíso inesperado para el entrenamiento de IA}
Japón podría convertirse en un nodo singular para el entrenamiento de IA a nivel mundial:
\begin{itemize}
    \item \textbf{Ley de derechos de autor}: permite explícitamente entrenar modelos con cualquier dato (una exención legal poco común)
    \item \textbf{Energía}: 9 GW de capacidad nuclear inactiva que puede reactivarse para centros de datos, una enorme fuente de electricidad barata
    \item \textbf{Hardware}: sin restricciones de importación de GPU (no está sujeto a los controles de exportación de Estados Unidos porque Japón es un país aliado)
    \item Varias empresas japonesas y capital extranjero ya están construyendo grandes instalaciones de entrenamiento de IA en Japón
\end{itemize}

Exención de derechos de autor + energía nuclear barata + GPU sin límite = un entorno de entrenamiento ideal. Es una oportunidad de arbitraje geográfico que casi nadie ha notado.
\end{knowledgebox}

También ocurren actividades de inteligencia dirigidas a investigadores de IA. Dylan señala, medio en broma: \textit{``como hombre soltero de veintitantos años... somos muy fáciles de corromper.''} Las operaciones de trampa (honey pot) dirigidas a jóvenes investigadores con conocimiento de frontera son una amenaza de seguridad real. Los servicios de inteligencia de distintos países compiten por el talento en IA, no para reclutar espías, sino para obtener secretos de entrenamiento y conocimientos arquitectónicos.

\subsection{La misión de código abierto de Nathan en AI2/Tulu}

El trabajo de Nathan en AI2 (Allen Institute for AI) representa el estándar más alto de IA de código abierto. Su proyecto Tulu fue el primero, entre los modelos que están en los 60 primeros lugares de Chatbot Arena, en publicar la receta \textbf{completa} de entrenamiento posterior (código + datos).

Resultado clave: Tulu aplica RLVR (RL con recompensas verificables) sobre el modelo base Llama, supera a la versión instruct oficial de Llama en razonamiento matemático y iguala el desempeño de DeepSeek V3. Esto demuestra que el \textbf{método de entrenamiento posterior en sí mismo} (y no solo la escala del modelo base) puede producir una diferencia enorme de desempeño.

OLMo es otro proyecto de AI2: un trabajo de preentrenamiento totalmente abierto (que complementa el entrenamiento posterior de Tulu). Con una escala de modelo mayor, es más fácil que el entrenamiento con RL despierte capacidades sólidas, que luego se pueden destilar a modelos pequeños.

La motivación central de Nathan: \textit{``No confío en la gente que dice 'confía en mí, hermano, vamos a hacer que la IA sea buena'.''} La IA es la tecnología más poderosa de nuestra vida: se necesita que más personas participen en darle forma, en lugar de confiar en que unas cuantas empresas decidan a puerta cerrada. La apertura ayuda a que investigadores fuera del campo de la IA, gobiernos y todos en general puedan entender lo que está ocurriendo.

\subsection{Resumen de la sección}

La economía de la industria de la IA se está reconfigurando a gran velocidad: la caída de costos de 1200 veces queda contrarrestada por la paradoja de Jevons, y la demanda total sigue creciendo. El contrabando de GPU se convirtió en una nueva industria subterránea. La destilación y la movilidad de talento hacen que los límites de la propiedad intelectual sean extremadamente difusos. La combinación singular de Japón de ley, energía y hardware crea una ventaja inesperada para el entrenamiento. El trabajo de Nathan en AI2/Tulu demuestra que el entrenamiento posterior totalmente abierto puede alcanzar el nivel de frontera.
